\documentclass[11pt]{article}
\usepackage[margin=1in]{geometry}
\usepackage{amsmath,amssymb}
\usepackage{graphicx}
\usepackage{booktabs}
\usepackage{hyperref}
\usepackage{xcolor}
\usepackage{caption}
\usepackage{subcaption}
\usepackage{multirow}

\title{When Do Transformers Learn Relational Structure?\\
A Synthetic Study of Affine Maps, Curriculum Design,\\
and Single-Exposure Learning}
\author{Hinrich Sch\"utze}
\date{May 2026}

\begin{document}
\maketitle

\begin{abstract}
We investigate when and how transformers learn structured relational
representations in synthetic settings with full experimental control.
The work spans two complementary lines of experiments.
In \textbf{Salzburg}, we test whether a single-layer transformer implements
relations as affine maps between subject and object embeddings, finding a
non-monotonic relationship between model capacity and structural solutions:
structure emerges only at an intermediate ``sweet spot.''
In \textbf{Graz}, we study single-exposure relational learning in a
compositional type system with 45 types and 2025 entities.  Through seven
experimental iterations (v1--v7), we establish that: (1) type-level
generalization is the dominant mechanism, accounting for over 90\% of relational
accuracy; (2) curriculum design --- deferring relation facts until entity
representations mature --- yields 50\% improvement over uniform training; and
(3) a single exposure to a relational fact adds only 3.9 percentage points
beyond what pure type-level generalization already provides.  The central
finding is that relational learning in transformers is overwhelmingly
\emph{type-level}: given sufficient training on a subset of entities, the model
generalizes the relational map to unseen entities at over 90\% accuracy, with
single exposures providing only per-subject refinement.
\end{abstract}

\tableofcontents
\newpage

%======================================================================
\section{Introduction}
%======================================================================

Large language models encode vast amounts of relational knowledge --- facts about
entities, their properties, and their relationships. But the nature of this encoding
remains debated: do models learn \emph{structured} representations that support
generalization, or do they \emph{memorize} specific associations?

Recent work has provided evidence for linear structure in LLM representations.
Hernandez et al.~\cite{hernandez2024} showed that certain relations are approximately
implemented as linear maps between subject and object embeddings. However, it remains
unclear (a) when such structure emerges, (b) what role training curriculum plays, and
(c) how much information a single exposure to a relational fact contributes beyond
type-level generalization.

We approach these questions through two synthetic studies:

\begin{itemize}
\item \textbf{Salzburg}: Can a transformer learn to implement relations as affine maps?
  We vary the hidden dimension and discover a non-monotonic capacity--structure
  relationship.
\item \textbf{Graz}: Can a transformer learn relational predictions from single
  exposures? We iterate through seven experimental designs (v1--v7), progressively
  isolating the mechanisms of relational learning and quantifying the contribution of
  curriculum design and single-exposure facts.
\end{itemize}

This report provides a comprehensive account of all experimental results, including
negative results and design iterations that informed the final findings.

%======================================================================
\section{Salzburg: Affine Relation Maps}
\label{sec:salzburg}
%======================================================================

\subsection{Setup}

The vocabulary consists of 100 subjects ($a_0, \ldots, a_{99}$), 100 objects
($b_0, \ldots, b_{99}$), 2 relation tokens ($r_0, r_1$), and a separator token ($x$).
Training sequences are 4-token triples $[x, r_k, a_i, b_j]$ where $b_j$ is the correct
object for subject $a_i$ under relation $r_k$.  Relation $r_0$ is the identity mapping
($a_i \mapsto b_i$) and $r_1$ is a cyclic shift by 50 ($a_i \mapsto b_{(i+50) \bmod 100}$).

The model is a single-layer transformer with 4 attention heads, trained with
next-token prediction (cross-entropy loss).  The hidden dimension $d$ varies
across experiments ($d \in \{8, 16, 32, 64\}$).

\paragraph{Evaluation.}
We measure: (1)~\emph{object prediction accuracy} --- does the model predict the
correct object? (2)~\emph{affine accuracy} --- we fit a per-relation affine map
$\hat{o} = W_k \cdot h_0(a_i) + b_k$ from input embeddings to output representations
and check whether the affine-predicted output falls in the correct Voronoi cell.

\subsection{Results: The Capacity--Structure Tradeoff}

All models achieve near-perfect object prediction ($\geq 99.5\%$) regardless of
dimension.  However, affine accuracy varies dramatically:

\begin{table}[h]
\centering
\caption{Salzburg dimension sweep. All models predict correctly, but affine structure
is non-monotonic in capacity.}
\label{tab:salzburg}
\begin{tabular}{lcccc}
\toprule
& dim=8 & dim=16 & dim=32 & dim=64 \\
\midrule
Parameters & 5,128 & 11,792 & 29,728 & 84,032 \\
Object accuracy & 99.5\% & 100\% & 100\% & 100\% \\
Affine accuracy & 73.0\% & 97.0\% & 100\% & 27.5\% \\
Existence test & 70.0\% & 99.0\% & 100\% & --- \\
\midrule
Regime & (A) & \multicolumn{2}{c}{Sweet spot} & (C) \\
\bottomrule
\end{tabular}
\end{table}

Three regimes are evident:
\begin{itemize}
\item \textbf{dim=8 --- No good solution exists (A).} 100 objects in 8 dimensions
  cannot be separated well enough for affine maps to work.
\item \textbf{dim=16--32 --- Sweet spot.} Affine maps are both representable and found
  by SGD, because they are the most parameter-efficient solution.
\item \textbf{dim=64 --- Memorization wins (C).} With 84K parameters for 200 facts
  (420 params/fact), the model memorizes without needing affine regularity.
\end{itemize}

\subsection{Auxiliary Loss Nudges Toward Structure}

At dim=16, adding a tiny auxiliary loss ($\lambda = 1.0$, stabilizing at ${\sim}0.02\%$
of the main loss) pushes affine accuracy from 97\% to \textbf{100\%}.  The architecture
\emph{can} implement the affine solution; standard SGD does not find it without explicit
pressure.

\subsection{Mechanistic Analysis}

A logit lens analysis reveals that attention routes information (copying subject and
relation embeddings) but cannot perform relation-dependent transformations. Only the
MLP's nonlinearity implements the gating needed to select which affine map to apply:

\begin{table}[h]
\centering
\caption{Logit lens: object prediction accuracy at intermediate stages (dim=16).}
\begin{tabular}{lcc}
\toprule
Stage & Accuracy & Affine $\cos$ \\
\midrule
$h_0$ (embedding only) & 17.2\% & --- \\
$h_0 + \text{attn}$ & 14.4\% & 0.987 \\
$h_0 + \text{attn} + \text{MLP}$ & 100\% & 0.773 \\
\bottomrule
\end{tabular}
\end{table}

The attention output is 0.987-affine in the input (mechanically linear), while the MLP
output is only 0.773-affine --- it provides the essential nonlinear gating.

%======================================================================
\section{Graz: Compositional Relational Learning}
\label{sec:graz}
%======================================================================

\subsection{Entity and Type Structure}

Entities are characterized by two independent attribute groups:
\begin{itemize}
\item \textbf{f-type}: a 2-element subset of 10 f-features $\to$ $\binom{10}{2} = 45$ types
\item \textbf{g-type}: a 2-element subset of 10 g-features $\to$ $\binom{10}{2} = 45$ types
\end{itemize}
Each (f-type, g-type) combination defines one entity type; with 45 entities per type,
there are $45 \times 45 = 2025$ entities total.

\paragraph{Relations.}
Two relations are defined via cyclic type shifts coprime with 45:
\begin{itemize}
\item $R_0$: shift the f-type index by 7 (mod 45)
\item $R_1$: shift the g-type index by 13 (mod 45)
\end{itemize}
Each subject maps to exactly 45 valid objects per relation (those sharing the target
type after the shift).  Chance level: $45/2025 = 2.22\%$.

\paragraph{Training sentences.}
Sequences of 48 tokens contain attribute sentences ($[x, e, a]$: predict attribute from
entity) and relation sentences ($[x, r_k, s_i, o_j]$: predict valid object from subject
and relation).

\paragraph{Model.}
Single-layer transformer, 4 attention heads, dim=64, 580,160 parameters.  All Graz
experiments (v1--v7) use this architecture.

\subsection{Evaluation}

We hold out 20\% of (subject, relation) pairs at random.  Held-out subjects still
appear in attribute sentences (their type is learned) but never in relation contexts.

\begin{itemize}
\item \textbf{Train relation accuracy}: top-1 accuracy on training (subj, rel) pairs
\item \textbf{Held-out relation accuracy}: top-1 accuracy on held-out pairs (main
  generalization metric)
\item \textbf{Affine probe}: post-hoc centroid probe trained on half the subjects,
  tested on the other half
\end{itemize}

%======================================================================
\section{Graz v1: Single-Exposure from Scratch (Negative Result)}
\label{sec:v1}
%======================================================================

\subsection{Design}

Phase~1 trains on attributes only for 50 epochs until 100\% accuracy.
Phase~2 introduces all relation sentences, each triple seen exactly once per epoch.
Setup: 6 features, 15 types, 225 entities, dim=16, 1 layer.

\subsection{Results}

\begin{table}[h]
\centering
\caption{v1: Total failure --- 0\% relation accuracy after 100 relation epochs.}
\begin{tabular}{rcccc}
\toprule
Rel Epoch & Loss & Train Rel & Held Rel & Attr Acc \\
\midrule
pre & --- & 0.000 & 0.000 & 1.000 \\
1 & 14.264 & 0.000 & 0.000 & 1.000 \\
10 & 12.577 & 0.000 & 0.000 & 0.106 \\
50 & 7.496 & 0.000 & 0.000 & 0.000 \\
100 & 7.079 & 0.000 & 0.000 & 0.000 \\
\bottomrule
\end{tabular}
\end{table}

\subsection{Diagnosis}

Two problems emerged:
\begin{enumerate}
\item \textbf{Relation token cold-start}: The relation tokens $r_0, r_1$ were never
  seen during phase~1. Their initial loss was 14.3 (vs.\ 2.5 for attribute sentences),
  creating extreme gradients that destroyed the learned representations.
\item \textbf{Catastrophic forgetting}: Attribute accuracy crashed from 100\% to 0\%
  and never recovered. The model lost all entity structure.
\end{enumerate}

\subsection{Lesson}

Relation tokens must be introduced alongside entity representations, not after.
Sequential phasing (attributes first, then relations) causes a distribution shift that
the small model cannot absorb.

%======================================================================
\section{Graz v2: Relation Token Warmup}
\label{sec:v2}
%======================================================================

\subsection{Design}

Phase~1 trains with a small ``pilot'' set of $k$ relation triples mixed in alongside
attributes. This teaches the model the relation sentence format without requiring
generalization.  Phase~2 then introduces all remaining triples.

Setup: 6 features, 15 types, 225 entities, dim=16, 360 training triples.
Chance level: 6.67\%.

\subsection{Phase 1: k-Triple Sweep}

\begin{table}[h]
\centering
\caption{Phase 1 results: all pilot triples memorized, no generalization until $k=32$.}
\begin{tabular}{rcccc}
\toprule
$k$ & Pilot Rel & Remaining Rel & Held Rel & Affine Test \\
\midrule
4 & 1.000 & 0.000 & 0.000 & 0.292 \\
8 & 1.000 & 0.000 & 0.000 & 0.296 \\
16 & 1.000 & 0.000 & 0.000 & 0.217 \\
32 & 1.000 & 0.067 & 0.067 & 0.230 \\
\bottomrule
\end{tabular}
\end{table}

\subsection{Phase 2: Single-Exposure Relation Epochs ($k=8$)}

After the $k=8$ warmup, remaining triples are presented once per epoch for 50 epochs:

\begin{table}[h]
\centering
\caption{v2 phase 2: slow learning, 16.7\% remaining accuracy after 50 epochs.}
\begin{tabular}{rccccc}
\toprule
Epoch & Loss & k-Rel & Remaining & Held & Attr \\
\midrule
pre & --- & 1.000 & 0.000 & 0.000 & 1.000 \\
2 & 5.541 & 1.000 & 0.056 & 0.022 & 1.000 \\
10 & 3.414 & 1.000 & 0.122 & 0.089 & 0.760 \\
50 & 2.978 & 1.000 & 0.167 & 0.078 & 0.954 \\
\bottomrule
\end{tabular}
\end{table}

\subsection{Findings}

The warmup solves the cold-start problem (initial loss drops from 14.3 to 5.95)
and prevents catastrophic forgetting.  However, with only 3 gradient steps per epoch
and 150 total steps over 50 epochs, the model barely exceeds chance on held-out triples.
The \textbf{gradient budget is the bottleneck}.

%======================================================================
\section{Graz v3: Large-Scale Joint Training}
\label{sec:v3}
%======================================================================

\subsection{Design}

The key insight from v2: the model needs far more gradient signal.  v3 scales to
$\binom{10}{2} = 45$ types, 2025 entities, and dim=64, with 145,800 training facts
(each seen exactly once) spread across 10 epochs.  Attributes and relations train
jointly from epoch~1.

\begin{table}[h]
\centering
\caption{v3 configuration vs.\ v2.}
\begin{tabular}{lcc}
\toprule
& v2 & v3 \\
\midrule
Entities & 225 & 2,025 \\
Types & 15 & 45 \\
Training facts & $\sim$360 & 145,800 \\
Parameters & 20,368 & 580,160 \\
Gradient steps & 150 & 9,760 \\
\bottomrule
\end{tabular}
\end{table}

\subsection{Results}

\begin{table}[h]
\centering
\caption{v3 epoch-by-epoch results (seed=42). Chance = 2.22\%.}
\label{tab:v3}
\begin{tabular}{rcccccc}
\toprule
Epoch & Loss & Attr & Train Rel & Held Rel & Affine (indiv.) & Affine (centroid) \\
\midrule
pre & --- & 0.003 & 0.010 & 0.014 & --- & --- \\
1 & 4.584 & 0.978 & 0.003 & 0.004 & 0.059 & --- \\
2 & 3.396 & 1.000 & 0.006 & 0.004 & 0.396 & --- \\
4 & 3.261 & 1.000 & 0.103 & 0.095 & 0.453 & --- \\
7 & 3.207 & 1.000 & 0.399 & 0.243 & 0.442 & --- \\
10 & 3.196 & 1.000 & 0.434 & 0.269 & 0.455 & 1.000 \\
\bottomrule
\end{tabular}
\end{table}

\subsection{Key Findings}

\paragraph{1. Genuine type-level generalization.}
Held-out accuracy reaches 26.9\% --- \textbf{12$\times$ chance}.  The model predicts
correct objects for (subject, relation) pairs it has never seen, purely from type
membership.

\paragraph{2. Perfect centroid probe.}
A type-centroid affine probe achieves 100\% on both train and test subjects.  The
model has learned a perfect linear map from each subject's hidden state to the centroid
of its 45 valid objects.  The gap between the centroid probe (100\%) and actual accuracy
(26.9--43.4\%) reflects the decoding bottleneck: the model points at the correct type
centroid, but the softmax over 2025 tokens must select one specific object from 45
candidates clustered around that centroid.

\paragraph{3. Phase transition at epoch 7.}
A sharp acceleration occurs between epochs 6--7: training accuracy jumps from 20.5\%
to 39.9\%.  After 6 epochs, the model has seen 60\% of training facts, providing enough
evidence for relational structure to consolidate.

\paragraph{4. The decoding bottleneck.}
Precision@45 analysis reveals the model finds only 12.5 of 45 valid objects in its
top-45 predictions. The bottleneck is not representation (centroid probe is perfect)
but output decoding: translating type-level knowledge into correct token predictions.

\subsection{Top-$k$ Coverage Analysis}

\begin{table}[h]
\centering
\caption{v3 top-$k$ analysis confirms perfect type clustering in both directions.}
\begin{tabular}{lcc}
\toprule
Metric & $k$ & Accuracy \\
\midrule
Entity $\to$ f-attributes & 2 & 1.000 \\
Entity $\to$ g-attributes & 2 & 1.000 \\
Entity $\to$ all attributes (joint) & 4 & 1.000 \\
Attribute $\to$ valid entities & 405 & 1.000 \\
\bottomrule
\end{tabular}
\end{table}

The model has learned complete type membership in both directions.  The relation
accuracy bottleneck lies entirely in relational decoding, not type clustering.

%======================================================================
\section{Graz v4: Reduced Object Pool (Negative Result)}
\label{sec:v4}
%======================================================================

\subsection{Hypothesis}

If the bottleneck is selecting among 45 valid objects out of 2025, reducing the object
pool should help.  We select a uniform subset of 675 objects (15 per f-type and 15 per
g-type, using K\"onig's theorem on a 45-regular bipartite graph).

\subsection{Results}

\begin{table}[h]
\centering
\caption{v4: reduced object pool. Accuracy is \emph{worse} than v3.}
\begin{tabular}{lcc}
\toprule
& v3 (2025 objects) & v4 (675 objects) \\
\midrule
Training facts & 145,800 & 48,600 \\
Train rel & 0.434 & 0.094 \\
Held-out rel & 0.269 & 0.080 \\
Chance & 2.22\% & 2.22\% \\
\bottomrule
\end{tabular}
\end{table}

\subsection{Lesson}

Reducing the object pool from 2025 to 675 also reduces training facts by $3\times$
(48,600 vs.\ 145,800).  The decreased relational signal far outweighs the easier
prediction problem.  The bottleneck is not object-pool size but \emph{gradient signal
from the relational facts themselves}.

%======================================================================
\section{Graz v5: Curriculum with Warmup}
\label{sec:v5}
%======================================================================

\subsection{Motivation}

The v3 phase transition at epoch~7 suggests that early-epoch relation facts are
partially wasted: attribute representations are still maturing in epochs 1--2, so
relation sentences presented then contribute less to relational learning.

\textbf{Key insight}: Single-exposure facts are a non-renewable resource.  Once
presented to immature representations, the learning opportunity is lost.

\subsection{Design}

\begin{itemize}
\item \textbf{Warmup (epochs 1--2)}: attributes + pilot set of 100 (subj, rel)
  pairs $\to$ 4,500 facts, repeated.  Purpose: learn relation format while
  representations mature.
\item \textbf{Main (epochs 3--10)}: all 145,800 single-exposure facts spread across
  8 epochs $\to$ 18,225 facts/epoch.
\end{itemize}

Same total compute as v3: 976 steps/epoch, 9,760 total steps.

\subsection{Results}

\begin{table}[h]
\centering
\caption{v5 epoch-by-epoch results. 29\% improvement over v3 with identical compute.}
\label{tab:v5}
\begin{tabular}{rccccc}
\toprule
Epoch & Phase & Loss & Attr & Train Rel & Held Rel \\
\midrule
pre & & --- & 0.003 & 0.010 & 0.014 \\
1 & warmup & 4.549 & 0.993 & 0.001 & 0.004 \\
2 & warmup & 3.361 & 1.000 & 0.000 & 0.000 \\
3 & main & 3.308 & 1.000 & 0.059 & 0.042 \\
5 & main & 3.240 & 1.000 & 0.243 & 0.157 \\
7 & main & 3.213 & 1.000 & 0.466 & 0.290 \\
10 & main & 3.201 & 1.000 & 0.560 & 0.347 \\
\bottomrule
\end{tabular}
\end{table}

\subsection{Key Finding: Scheduling Matters as Much as Quantity}

\begin{table}[h]
\centering
\caption{v5 vs.\ v3: same data, same compute, different scheduling.}
\begin{tabular}{lccc}
\toprule
& v3 & v5 & Improvement \\
\midrule
Train rel & 0.434 & 0.560 & +29\% \\
Held-out rel & 0.269 & 0.347 & +29\% \\
Held-out / chance & 12.1$\times$ & 15.6$\times$ & \\
\bottomrule
\end{tabular}
\end{table}

The only difference between v3 and v5 is \emph{when} facts are presented.  Deferring
single-exposure relation facts to epochs with mature entity representations yields a
29\% improvement in both training and held-out accuracy.  After 2 warmup epochs, the
pilot set shows 0\% relation accuracy --- confirming that relation facts presented to
immature representations are indeed wasted.

%======================================================================
\section{Graz v6: Graduated Curriculum with Pilot Scaffold}
\label{sec:v6}
%======================================================================

\subsection{Motivation}

v5's 2-epoch warmup establishes entity representations, but the model enters the main
phase with no relational knowledge.  What if we extend the pilot phase to actually
\emph{learn} the relational pattern before exposing single-exposure facts?

\subsection{Design}

\begin{itemize}
\item \textbf{Pilot phase (epochs 1--5)}: 100 (subj, rel) pairs $\to$ 4,500 facts,
  repeated each epoch.  Goal: learn the relational map from a small scaffold.
\item \textbf{Main phase (epochs 6--13)}: all 145,800 single-exposure facts across
  8 epochs $\to$ 18,225 facts/epoch.
\end{itemize}

Total: 13 epochs, 12,688 steps (30\% more compute than v3/v5).

\subsection{Results}

\begin{table}[h]
\centering
\caption{v6: pilot scaffold enables dramatically better single-exposure learning.}
\label{tab:v6}
\begin{tabular}{rcccccc}
\toprule
Epoch & Phase & Attr & Pilot & Non-pilot & Held & Affine \\
\midrule
1 & pilot & 0.994 & 0.010 & 0.001 & 0.004 & 0.099 \\
2 & pilot & 1.000 & 0.000 & 0.000 & 0.000 & 0.457 \\
3 & pilot & 1.000 & 0.050 & 0.000 & 0.000 & 0.492 \\
5 & pilot & 1.000 & 0.620 & 0.000 & 0.000 & 0.430 \\
\midrule
6 & main & 1.000 & 0.750 & 0.185 & 0.141 & 0.420 \\
8 & main & 1.000 & 0.900 & 0.583 & 0.348 & 0.430 \\
11 & main & 1.000 & 0.920 & 0.691 & 0.405 & 0.436 \\
13 & main & 1.000 & 0.920 & 0.706 & 0.393 & 0.441 \\
\bottomrule
\end{tabular}
\end{table}

\subsection{Key Findings}

\paragraph{1. The scaffold builds a reusable relational prior.}
During the pilot phase, 100 triples reach 62\% accuracy while non-pilot accuracy
stays at exactly 0\%.  No leakage --- the model learns subject-specific relational
mappings, not just type patterns.

\paragraph{2. Immediate transfer at main phase onset.}
At epoch~6 (first main epoch), non-pilot accuracy jumps from 0\% to 18.5\%.  These
triples have never been seen.  The 18.5\% first-epoch transfer (vs.\ 2.22\% chance)
is strong evidence that the pilot scaffold teaches the type-level relational map.

\paragraph{3. Pilot accuracy rises during main phase.}
Pilot triples improve from 62\% to 92\% during the main phase, even though no new
pilot facts are presented.  Additional facts about different subjects of the same
types reinforce the pilot's relational pattern --- confirming type-level learning.

\paragraph{4. v3 $\to$ v5 $\to$ v6 progression.}

\begin{table}[h]
\centering
\caption{Progressive improvement through curriculum refinement.}
\begin{tabular}{llccccc}
\toprule
Version & Strategy & Train Rel & Held Rel & Epochs & Steps \\
\midrule
v3 & Uniform exposure & 0.434 & 0.269 & 10 & 9,760 \\
v5 & 2-epoch warmup & 0.560 & 0.347 & 10 & 9,760 \\
v6 & 5-epoch scaffold & 0.706$^*$ & 0.405 & 13 & 12,688 \\
\bottomrule
\multicolumn{6}{l}{\small $^*$Non-pilot accuracy (single-exposure triples only).}
\end{tabular}
\end{table}

Each curriculum refinement improves held-out accuracy: +29\% (v5 vs.\ v3), +17\%
(v6 vs.\ v5), +50\% total (v6 vs.\ v3).

%======================================================================
\section{Graz v7: Three-Way Split --- Isolating Single-Exposure Value}
\label{sec:v7}
%======================================================================

\subsection{Motivation}

The v3--v6 experiments conflate two sources of relational learning: (a) the type-level
map learned from repeated facts, and (b) per-subject information from single exposures.
v7 is designed to cleanly separate these by introducing a three-way split:

\begin{itemize}
\item \textbf{Train (60\%)}: 2,430 (subj, rel) pairs, repeated over 10 epochs
\item \textbf{Single-exposure (20\%)}: 810 pairs $\to$ 36,450 facts, each seen once
\item \textbf{Held-out (20\%)}: 810 pairs, never seen
\end{itemize}

The key comparison: after Phase~2 (single-exposure), how much does the single-exposure
group improve \emph{relative to held-out} (which received no exposure)?

\subsection{Design}

\begin{itemize}
\item \textbf{Phase 1 (10 epochs)}: Attributes + 60\% train relations, repeated.
  Cosine LR from 3e-4, 200-step warmup.  $p_\text{relation} = 0.20$.
\item \textbf{Phase 2 (1 epoch)}: 36,450 single-exposure facts, each seen exactly
  once.  Constant LR = 1e-4.
\end{itemize}

Total: 10,736 steps (976/epoch).

\subsection{Results}

\begin{table}[h]
\centering
\caption{v7 Phase 1: learning dynamics. Single and held-out are identical (neither
seen), confirming no information leakage.}
\label{tab:v7p1}
\begin{tabular}{rcccccc}
\toprule
Epoch & Loss & Attr & Train 60\% & Single 20\% & Held 20\% \\
\midrule
pre & --- & 0.003 & 0.010 & 0.011 & 0.014 \\
1 & 4.766 & 0.416 & 0.000 & 0.000 & 0.000 \\
2 & 3.675 & 1.000 & 0.049 & 0.027 & 0.031 \\
3 & 3.481 & 1.000 & 0.630 & 0.326 & 0.337 \\
4 & 3.359 & 1.000 & 0.908 & 0.616 & 0.662 \\
5 & 3.291 & 1.000 & 0.994 & 0.801 & 0.821 \\
7 & 3.242 & 1.000 & 1.000 & 0.861 & 0.873 \\
10 & 3.227 & 1.000 & 1.000 & 0.906 & 0.906 \\
\bottomrule
\end{tabular}
\end{table}

\begin{table}[h]
\centering
\caption{v7 Phase 2: the value of a single exposure.}
\label{tab:v7p2}
\begin{tabular}{lccc}
\toprule
& Train 60\% & Single 20\% & Held 20\% \\
\midrule
After Phase 1 & 1.000 & 0.906 & 0.906 \\
After Phase 2 & 0.999 & \textbf{0.995} & \textbf{0.956} \\
Change & $-$0.001 & +0.089 & +0.049 \\
\bottomrule
\end{tabular}
\end{table}

\subsection{Key Findings}

\paragraph{1. Type-level generalization reaches 90.6\% without any single-exposure facts.}
After Phase~1, both the single-exposure and held-out groups score 90.6\% (identical,
as expected).  The model learns the relational map from 60\% of entities well enough
to generalize to the remaining 40\% at over 90\% accuracy ---
\textbf{40.8$\times$ chance}.

This is dramatically higher than v3--v6 held-out accuracy (27--41\%).  The difference:
Phase~1 trains on 60\% of relations with full repetition (10 epochs), giving the model
sufficient signal to learn the complete type-level map.

\paragraph{2. A single exposure adds 8.9 percentage points.}
Single-exposure facts go from 90.6\% $\to$ 99.5\% after one pass through 36,450 facts.
This +8.9pp gain represents per-subject information that the type-level map cannot
capture.

\paragraph{3. Spillover benefits held-out by 4.9pp.}
Held-out accuracy rises from 90.6\% $\to$ 95.6\% despite those facts never being seen.
The additional 36,450 facts in Phase~2 strengthen the type-level map itself.

\paragraph{4. Net single-exposure advantage: 3.9 percentage points.}

\begin{table}[h]
\centering
\caption{The net value of a single exposure, controlling for spillover.}
\begin{tabular}{lcc}
\toprule
Group & After Phase 2 & Source of knowledge \\
\midrule
Single (seen once) & 99.5\% & Type-level + per-subject detail \\
Held (never seen) & 95.6\% & Type-level only \\
\textbf{Difference} & \textbf{3.9pp} & Per-subject information from single exposure \\
\bottomrule
\end{tabular}
\end{table}

\paragraph{5. Phase 1 learning dynamics.}
The train/generalization ratio converges from 1.90 (epoch~3) to 1.10 (epoch~10),
indicating that the model's relational knowledge becomes increasingly type-level as
training progresses.  By epoch~10, 91\% of what the model knows about trained relations
transfers to unseen subjects.

%======================================================================
\section{Summary of Graz Progression}
\label{sec:summary}
%======================================================================

\begin{table}[h]
\centering
\caption{Complete Graz progression from v1 (failure) to v7 (clean isolation).}
\label{tab:all_versions}
\begin{tabular}{llccccl}
\toprule
Ver. & Strategy & Entities & Train & Held & Steps & Key insight \\
\midrule
v1 & Sequential phases & 225 & 0.000 & 0.000 & 300 & Cold-start kills \\
v2 & Warmup + exposure & 225 & 0.167 & 0.078 & 150 & Gradient-limited \\
v3 & Joint training & 2,025 & 0.434 & 0.269 & 9,760 & Scale works \\
v4 & Reduced pool & 2,025 & 0.094 & 0.080 & 9,760 & Less signal hurts \\
v5 & Curriculum & 2,025 & 0.560 & 0.347 & 9,760 & Timing matters \\
v6 & Scaffold & 2,025 & 0.706 & 0.405 & 12,688 & Prior helps \\
v7 & 3-way split & 2,025 & 1.000 & 0.956 & 10,736 & Type $\gg$ instance \\
\bottomrule
\end{tabular}
\end{table}

%======================================================================
\section{Discussion}
\label{sec:discussion}
%======================================================================

\subsection{The Central Finding: Relational Learning is Type-Level}

The v7 experiment provides the cleanest answer to the original research question.
A transformer trained on 60\% of (subject, relation) pairs with repeated exposure
generalizes to the remaining 40\% at 90.6\% accuracy.  A single additional exposure
adds only 3.9pp of per-subject information beyond this.

This reframes ``single-exposure relational learning'': the model does not learn the
relation from a single exposure.  It already knows the relation from the type structure.
The single exposure merely fills in subject-specific residual information within a
well-learned relational scaffold.

\subsection{Curriculum Design Matters}

The progression from v3 to v6 demonstrates that the scheduling of learning signal
matters as much as its quantity:
\begin{itemize}
\item v5 (+29\% over v3): same data, same compute, deferred scheduling
\item v6 (+50\% over v3): 30\% more compute, pilot scaffold
\end{itemize}
The key principle: \textbf{single-exposure facts are a non-renewable resource}.
Presenting them to immature representations wastes the learning opportunity.  Optimal
curriculum design ensures entity representations and relational priors are established
before single-exposure facts are consumed.

\subsection{The Capacity--Structure Connection}

The Salzburg results (Section~\ref{sec:salzburg}) and the early Graz experiments
(not reported in detail here; see Appendix) show the same non-monotonic pattern:
structured solutions emerge only when model capacity is matched to task complexity.
However, v7 reveals that once sufficient training data is provided, the
capacity concern becomes secondary: the dim=64 model (which in the dimension-sweep
experiments memorized rather than generalized) achieves 90.6\% type-level
generalization when trained on enough relations.  The capacity--structure tradeoff
is thus modulated by data quantity.

\subsection{Three Mechanisms of Relational Knowledge}

v7 cleanly separates three mechanisms:
\begin{enumerate}
\item \textbf{Type-level generalization (90.6\%)}:  From repeated training on 60\% of
  entities, the model learns a complete type-to-type relational map that transfers to
  unseen entities.
\item \textbf{Spillover reinforcement (+4.9pp)}:  Additional single-exposure facts
  strengthen the type-level map even for entities that never appear in those facts.
\item \textbf{Per-subject detail (+3.9pp)}:  A single exposure provides subject-specific
  information that pure generalization cannot supply.
\end{enumerate}

\subsection{Implications}

\paragraph{For understanding LLMs.}
Real LLMs encounter most relational facts only a few times in their training corpus.
Our results suggest that for relations with clear type structure (``capital of,''
``CEO of''), the model may learn the relation primarily from type-level patterns
(countries map to cities, companies map to people) rather than from individual fact
exposures.  The individual facts then serve primarily as per-entity refinements.

\paragraph{For knowledge editing.}
If relational knowledge is primarily type-level, then editing a single fact (e.g.,
changing a capital city) requires overcoming the type-level prior.  This may explain
why knowledge editing techniques often exhibit ``ripple effects'' --- the type-level
structure resists localized changes.

\paragraph{For data efficiency.}
The curriculum results suggest that fact presentation order matters for data-efficient
learning.  Training on type-representative examples first (establishing the relational
scaffold) and then presenting the long tail of specific facts may yield better
generalization than uniform random training.

\subsection{Limitations}

\begin{itemize}
\item All experiments use a single-layer transformer.  Multi-layer models may exhibit
  different dynamics.
\item The synthetic type system has perfectly clean structure.  Real-world relations
  are noisier and may not decompose as cleanly into type-level and instance-level
  components.
\item The 45-type system with 45 entities per type creates a balanced structure.
  Skewed distributions (few entities per type, or many types with few instances)
  might change the balance between type-level and per-subject learning.
\item We use dim=64 throughout the v3--v7 experiments.  The Salzburg results suggest
  this may be in the ``memorization regime'' for the capacity sweep, but the v7
  result shows that with enough data, generalization succeeds regardless.
\end{itemize}

%======================================================================
\section{Conclusion}
\label{sec:conclusion}
%======================================================================

Through seven experimental iterations in the Graz setting, we have established that:

\begin{enumerate}
\item \textbf{Relational learning is overwhelmingly type-level.}  A model trained on
  60\% of entities generalizes to unseen entities at 90.6\% accuracy.  Per-subject
  information from single exposures adds only 3.9pp.

\item \textbf{Curriculum design yields large gains.}  Deferring single-exposure facts
  until entity representations mature improves generalization by 50\% (v6 vs.\ v3)
  with modest additional compute.

\item \textbf{The relational map is a reusable scaffold.}  Once learned from a small
  pilot set, it enables efficient absorption of new facts: non-pilot triples reach
  70.6\% after a single exposure, and a full 60\% training set enables 90.6\%
  pure generalization.

\item \textbf{Structure vs.\ memorization depends on capacity-to-data ratio.}  The
  Salzburg results confirm that structured solutions emerge only at intermediate
  capacity.  But with sufficient relational data (v7), even high-capacity models
  learn generalizable type-level structure.
\end{enumerate}

The central message: transformers learn relations primarily as type-level maps.
Individual fact exposure serves as refinement, not as the primary learning mechanism.

\appendix

%======================================================================
\section{Graz Phase 3 Experiments: Feature Reduction and Contrastive Learning}
\label{app:progression}
%======================================================================

Prior to the one-epoch experiments (v1--v7), we conducted a series of experiments
(``Phase 3'') exploring the transition from memorization to generalization in a
smaller setting.  These experiments motivated the larger-scale one-epoch line of work.

\subsection{Progression from Memorization to Generalization}

\begin{table}[h]
\centering
\caption{Phase 3 experimental progression.}
\begin{tabular}{llccccc}
\toprule
Phase & Key change & Features & Entities & dim & Held-out & $\times$ chance \\
\midrule
Phase 2 & Subject holdout & 10 & 120 & 32 & 0\% & --- \\
Phase 3 & Triple holdout & 10 & 120 & 32 & 35\% & 42$\times$ \\
Phase 3j & 7 features, $k{=}2$ & 7 & 441 & 32 & 50.6\% & 10.6$\times$ \\
Phase 3l & + contrastive & 7 & 441 & 64 & 69.3\% & 14.6$\times$ \\
Phase 3n & 6 features & 6 & 225 & 64 & 87.8\% & 13.2$\times$ \\
Phase 3o3 & 5 features, dim=16 & 5 & 100 & 16 & 90.0\% & 9.0$\times$ \\
\bottomrule
\end{tabular}
\end{table}

\subsection{Dimension Sweep (Phase 3o)}

The Graz dimension sweep at 5 features / 100 entities shows the same
non-monotonic capacity--structure pattern as Salzburg:

\begin{table}[h]
\centering
\caption{Graz dimension sweep (5 features, 100 entities). Chance = 10\%.}
\begin{tabular}{lcccc}
\toprule
& dim=8 & dim=16 & dim=32 & dim=64 \\
\midrule
Held-out accuracy & 42.5\% & \textbf{90.0\%} & 85.0\% & 67.5\% \\
Contrastive $+$/$-$ & 85.2/94.0 & 88.8/97.8 & 88.8/99.2 & 90.0/97.6 \\
Affine probe (test) & 26.0\% & 50.0\% & 55.0\% & 63.0\% \\
Partner similarity & +0.22 & +0.43 & +0.18 & +0.10 \\
\bottomrule
\end{tabular}
\end{table}

\subsection{Contrastive Learning Ablation}

Contrastive sentences (``is object $b_j$ related to subject $a_i$ via $r_k$?  yes/no'')
are critical for generalization in the small-scale setting:

\begin{table}[h]
\centering
\caption{Effect of contrastive learning (7 features, 441 entities). Chance = 4.8\%.}
\begin{tabular}{lccc}
\toprule
& No contrastive & Contrastive, 6.25K seq & Contrastive, 62.5K seq \\
\midrule
Held-out accuracy & 50.6\% & 30.7\% & \textbf{69.3\%} \\
$\times$ chance & 10.6$\times$ & 6.4$\times$ & 14.6$\times$ \\
\bottomrule
\end{tabular}
\end{table}

With insufficient data, contrastive sentences \emph{hurt} (competing for capacity).
With 10$\times$ more data, they drive strong compositional generalization.

%======================================================================
\section{Salzburg: Additional Mechanistic Details}
\label{app:salzburg}
%======================================================================

\subsection{Sublayer Decomposition}

\begin{table}[h]
\centering
\caption{Affine predictability of sublayer components from input embedding (dim=16).}
\begin{tabular}{lcc}
\toprule
Target & $\cos(\hat{o}, o)$ & Interpretation \\
\midrule
Attention output & 0.987 & Near-perfectly affine \\
MLP output & 0.773 & Nonlinear gating \\
Full residual & 0.735 & Sum of all components \\
\bottomrule
\end{tabular}
\end{table}

Attention is mechanically linear (weighted sum of value-projected inputs).  The MLP's
GELU nonlinearity implements the relation-dependent gating: ``if relation = $k$,
apply transformation $W_k$.''  This explains why post-attention accuracy (14.4\%)
is no better than chance despite the attention output being highly structured.

\begin{thebibliography}{9}

\bibitem{hernandez2024}
E.~Hernandez, A.~S.~Sharma, T.~Haklay, K.~Meng, M.~Wattenberg, J.~Andreas,
Y.~Belinkov, and D.~Bau.
\newblock Linearity of relation decoding in transformer language models.
\newblock In \emph{ICLR}, 2024.

\bibitem{power2022}
A.~Power, Y.~Burda, H.~Edwards, I.~Babuschkin, and V.~Misra.
\newblock Grokking: Generalization beyond overfitting on small algorithmic datasets.
\newblock \emph{arXiv preprint arXiv:2201.02177}, 2022.

\bibitem{lake2018}
B.~M.~Lake and M.~Baroni.
\newblock Generalization without systematicity: On the compositional skills of
  sequence-to-sequence recurrent networks.
\newblock In \emph{ICML}, 2018.

\bibitem{kim2020}
N.~Kim and A.~Linzen.
\newblock COGS: A compositional generalization challenge based on semantic
  interpretation.
\newblock In \emph{EMNLP}, 2020.

\bibitem{tishby2000}
N.~Tishby, F.~C.~Pereira, and W.~Bialek.
\newblock The information bottleneck method.
\newblock \emph{arXiv preprint physics/0004057}, 2000.

\bibitem{kandpal2023}
N.~Kandpal, H.~Deng, A.~Roberts, E.~Wallace, and C.~Raffel.
\newblock Large language models struggle to learn long-tail knowledge.
\newblock In \emph{ICML}, 2023.

\end{thebibliography}

\end{document}
